% Generated by publication/build_sections.py; do not edit by hand.
\paragraph{Informal verdict.}
\textbf{A --- Complete solution}, with \textbf{high confidence}.

\paragraph{Lean coverage.}
Lean unconditionally proves existence and uniqueness of the universal integer polynomial for every n, m, and i. It also kernel-checks, for every positive m and primitive m-th root, the finite root-of-unity coefficient and expanded-sum formulas, which apply over the complex numbers for every canonical m >= 2, plus the complete m = 2 family, boundary cases, Frobenius specialization, and worked examples. The unformalized companion/resultant and Newton/partition presentations are alternative descriptions, not additional requirements of Q831.

\subsection{Audited informal answer}
Q831 --- Complete solution
\begin{enumerate}
\item Faithful restatement
\end{enumerate}
Let \(\Sigma_1,\ldots,\Sigma_n\) be the elementary symmetric polynomials in \(X_1,\ldots,X_n\). For each integer \(m\ge 2\), express explicitly

\[
\Sigma_i(X_1^m,X_2^m,\ldots,X_n^m)
\]

in terms of the original elementary symmetric polynomials
\(\Sigma_1(X),\ldots,\Sigma_n(X)\).
The PDF contains no subparts or additional hypotheses. There is no apparent typo. The only minor ambiguity is the word ``explicit''; below it is answered in three strong senses: a closed coefficient formula, an integral determinant/resultant formula, and a finite partition formula. Luc Pommeret

\begin{enumerate}
\item Notation and main theorem
\end{enumerate}
Write

\[
e_r=\Sigma_r(X_1,\ldots,X_n),\qquad e_0=1,
\]

and set \(e_r=0\) when \(r<0\) or \(r>n\). Define

\[
e_i^{[m]}:=\Sigma_i(X_1^m,\ldots,X_n^m).
\]

Thus

\[
e_i^{[m]}
=
\sum_{1\le j_1<\cdots<j_i\le n}
(X_{j_1}\cdots X_{j_i})^m.
\]

Theorem
For every \(n\ge1\), \(m\ge1\), and \(0\le i\le n\), there is a unique polynomial

\[
\Phi_{m,i}^{(n)}\in\mathbb Z[Y_1,\ldots,Y_n]
\]

such that

\[
e_i^{[m]}=\Phi_{m,i}^{(n)}(e_1,\ldots,e_n).
\]

It has the following explicit descriptions.
2.1 Root-of-unity formula
Let \(\zeta\) be any primitive \(m\)-th root of unity, and put \(Y_0=1\). Then

\[
\boxed{
\Phi_{m,i}^{(n)}(Y_1,\ldots,Y_n)
=
(-1)^{(m+1)i}
\sum_{\substack{0\le k_0,\ldots,k_{m-1}\le n\\
k_0+\cdots+k_{m-1}=mi}}
\zeta^{\,\sum_{r=0}^{m-1}r k_r}
\prod_{r=0}^{m-1}Y_{k_r}.
}
\tag{2.1}
\]

Although roots of unity occur in the displayed expression, all its coefficients are rational integers; it is independent of the chosen primitive root \(\zeta\).
Equivalently, if

\[
E(z)=1+Y_1z+\cdots+Y_nz^n,
\]

then

\[
\boxed{
\Phi_{m,i}^{(n)}
=
(-1)^{(m+1)i}
[z^{mi}]
\prod_{r=0}^{m-1}E(\zeta^rz).
}
\tag{2.2}
\]

2.2 Companion-matrix and resultant formulas
Let

\[
P_Y(Z)
=
Z^n-Y_1Z^{n-1}+Y_2Z^{n-2}-\cdots+(-1)^nY_n,
\]

and let \(C=C(Y)\) be its companion matrix:

\[
C=
\begin{pmatrix}
0&0&\cdots&0&(-1)^{n+1}Y_n\\
1&0&\cdots&0&(-1)^nY_{n-1}\\
0&1&\cdots&0&(-1)^{n-1}Y_{n-2}\\
\vdots&&\ddots&\vdots&\vdots\\
0&0&\cdots&1&Y_1
\end{pmatrix}.
\]

Then

\[
\boxed{
\det(TI_n-C^m)
=
\sum_{i=0}^n(-1)^i\Phi_{m,i}^{(n)}T^{\,n-i}.
}
\tag{2.3}
\]

Consequently,

\[
\boxed{
\Phi_{m,i}^{(n)}
=
(-1)^i[T^{\,n-i}]\det(TI_n-C^m).
}
\tag{2.4}
\]

Equivalently, with the convention that the resultant of a monic polynomial \(P\) with \(G\) is the determinant of multiplication by \(G\) on the quotient by \(P\),

\[
\boxed{
\operatorname{Res}_Z\!\left(
Z^n-Y_1Z^{n-1}+\cdots+(-1)^nY_n,\,
T-Z^m
\right)
=
\sum_{i=0}^n(-1)^i\Phi_{m,i}^{(n)}T^{\,n-i}.
}
\tag{2.5}
\]

This is a single integral determinant and involves neither division nor roots of unity.

\begin{enumerate}
\item Proof of the theorem
\end{enumerate}
3.1 The companion-matrix identity
Work first in the universal polynomial ring

\[
A=\mathbb Z[Y_1,\ldots,Y_n].
\]

The matrix \(C\) is the matrix, in the basis

\[
1,Z,\ldots,Z^{n-1},
\]

of multiplication by \(Z\) in the free \(A\)-module

\[
A[Z]/(P_Y).
\]

Indeed, multiplication by \(Z\) shifts the first \(n-1\) basis vectors, while the relation

\[
Z^n
=
Y_1Z^{n-1}-Y_2Z^{n-2}
+\cdots+(-1)^{n+1}Y_n
\]

gives precisely the last column of \(C\).
Now specialize

\[
Y_r=e_r(X_1,\ldots,X_n).
\]

Then

\[
P_Y(Z)
=
\prod_{j=1}^n(Z-X_j).
\]

Over the field \(K=\mathbb Q(X_1,\ldots,X_n)\), the elements \(X_1,\ldots,X_n\) are pairwise distinct, and the Chinese remainder theorem gives

\[
K[Z]/(P_Y)
\cong
\prod_{j=1}^nK,
\qquad
f(Z)\longmapsto
\bigl(f(X_1),\ldots,f(X_n)\bigr).
\]

Under this isomorphism, multiplication by \(Z\) becomes multiplication by the diagonal matrix

\[
\operatorname{diag}(X_1,\ldots,X_n).
\]

Therefore multiplication by \(Z^m\) has eigenvalues

\[
X_1^m,\ldots,X_n^m,
\]

and hence

\[
\det(TI_n-C^m)
=
\prod_{j=1}^n(T-X_j^m).
\]

Expanding the right-hand side gives

\[
\prod_{j=1}^n(T-X_j^m)
=
\sum_{i=0}^n
(-1)^i
\Sigma_i(X_1^m,\ldots,X_n^m)T^{n-i}.
\]

Comparing coefficients proves (2.3)--(2.5), and in particular proves the existence of integral polynomials \(\Phi_{m,i}^{(n)}\).
3.2 Uniqueness
The elementary symmetric polynomials \(e_1,\ldots,e_n\) are algebraically independent over \(\mathbb Z\). Here is a direct verification.
Use the lexicographic order \(X_1>\cdots>X_n\). Then

\[
\operatorname{LM}(e_r)=X_1X_2\cdots X_r.
\]

Consequently, for \(\alpha=(\alpha_1,\ldots,\alpha_n)\),

\[
\operatorname{LM}
\left(
e_1^{\alpha_1}\cdots e_n^{\alpha_n}
\right)
=
X_1^{b_1}\cdots X_n^{b_n},
\qquad
b_j=\sum_{r=j}^n\alpha_r.
\]

The vector \((b_1,\ldots,b_n)\) determines \(\alpha\), since

\[
\alpha_j=b_j-b_{j+1},
\qquad b_{n+1}=0.
\]

Thus distinct monomials in the \(e_j\)'s have distinct leading monomials after substitution, so no nonzero polynomial in \(Y_1,\ldots,Y_n\) can vanish after substituting \(Y_j=e_j\). This proves uniqueness.
3.3 Derivation of the root-of-unity formula
We have

\[
E(z):=\sum_{k=0}^ne_kz^k
=
\prod_{j=1}^n(1+X_jz).
\]

For a primitive \(m\)-th root \(\zeta\),

\[
\prod_{r=0}^{m-1}(1-\zeta^ru)=1-u^m.
\]

Taking \(u=-zX_j\), we obtain

\[
\prod_{r=0}^{m-1}(1+\zeta^rzX_j)
=
1-(-zX_j)^m
=
1+(-1)^{m+1}z^mX_j^m.
\]

Therefore

\[
\begin{aligned}
\prod_{r=0}^{m-1}E(\zeta^rz)
&=
\prod_{r=0}^{m-1}\prod_{j=1}^n
(1+\zeta^rzX_j)\\
&=
\prod_{j=1}^n
\left(1+(-1)^{m+1}z^mX_j^m\right)\\
&=
\sum_{i=0}^n
(-1)^{(m+1)i}
e_i^{[m]}z^{mi}.
\end{aligned}
\]

Hence

\[
e_i^{[m]}
=
(-1)^{(m+1)i}
[z^{mi}]
\prod_{r=0}^{m-1}E(\zeta^rz),
\]

which is (2.2).
Expanding each factor,

\[
E(\zeta^rz)
=
\sum_{k=0}^ne_k\zeta^{rk}z^k,
\]

gives (2.1).
The left side is already known from the companion formula to be an integral polynomial in the \(e_j\)'s. By algebraic independence, the root-of-unity expression is therefore precisely the same integral polynomial \(\Phi_{m,i}^{(n)}\). This also proves that it is independent of the chosen primitive root.

\begin{enumerate}
\item A formula involving only integer sums and partitions
\end{enumerate}
The determinant formula is division-free. There is also a fully expanded formula through the power sums.
Let

\[
p_q=X_1^q+\cdots+X_n^q.
\]

For \(q\ge1\), define the Newton polynomial

\[
\boxed{
N_q(Y)
=
\sum_{\substack{a_1,\ldots,a_q\ge0\\
a_1+2a_2+\cdots+qa_q=q}}
(-1)^{q-A}
\frac{q(A-1)!}{a_1!\cdots a_q!}
Y_1^{a_1}\cdots Y_q^{a_q},
}
\tag{4.1}
\]

where

\[
A=a_1+\cdots+a_q,
\qquad
Y_r=0\quad(r>n).
\]

Then

\[
p_q=N_q(e_1,\ldots,e_n).
\]

The apparently fractional coefficient in (4.1) is an integer, because

\[
\frac{q(A-1)!}{a_1!\cdots a_q!}
=
\sum_{\substack{1\le s\le q\\a_s>0}}
s\,
\frac{(A-1)!}
{a_1!\cdots(a_s-1)!\cdots a_q!}.
\]

The first Newton polynomials are

\[
\begin{aligned}
N_1&=Y_1,\\
N_2&=Y_1^2-2Y_2,\\
N_3&=Y_1^3-3Y_1Y_2+3Y_3,\\
N_4&=Y_1^4-4Y_1^2Y_2+2Y_2^2+4Y_1Y_3-4Y_4.
\end{aligned}
\]

Now let \(\lambda\vdash i\) be a partition of \(i\). If \(m_r(\lambda)\) denotes the multiplicity of \(r\) in \(\lambda\), put

\[
z_\lambda=\prod_{r\ge1}r^{m_r(\lambda)}m_r(\lambda)!.
\]

Then

\[
\boxed{
\Phi_{m,i}^{(n)}
=
\sum_{\lambda\vdash i}
\frac{(-1)^{i-\ell(\lambda)}}{z_\lambda}
\prod_{a=1}^{\ell(\lambda)}
N_{m\lambda_a}.
}
\tag{4.2}
\]

Although rational coefficients occur termwise in (4.2), their sum is the integral polynomial given by the companion determinant.
Proof
All power-series calculations below are formal, in
\(\mathbb Q[X_1,\ldots,X_n][[u]]\); no convergence is involved.
First,

\[
\begin{aligned}
\log\left(\sum_{r=0}^ne_ru^r\right)
&=
\sum_{j=1}^n\log(1+X_ju)\\
&=
\sum_{q\ge1}
\frac{(-1)^{q-1}}q\,p_qu^q.
\end{aligned}
\tag{4.3}
\]

On the other hand,

\[
\log\left(1+\sum_{r=1}^ne_ru^r\right)
=
\sum_{A\ge1}
\frac{(-1)^{A-1}}A
\left(\sum_{r=1}^ne_ru^r\right)^A.
\]

Extracting the coefficient of \(u^q\) gives exactly (4.1).
Next,

\[
\begin{aligned}
\sum_{i=0}^ne_i^{[m]}t^i
&=
\prod_{j=1}^n(1+X_j^mt)\\
&=
\exp\left(
\sum_{r\ge1}
\frac{(-1)^{r-1}}r\,p_{mr}t^r
\right).
\end{aligned}
\tag{4.4}
\]

Expanding the exponential and collecting terms according to partitions of \(i\) gives

\[
e_i^{[m]}
=
\sum_{\lambda\vdash i}
\frac{(-1)^{i-\ell(\lambda)}}{z_\lambda}
\prod_{a=1}^{\ell(\lambda)}p_{m\lambda_a}.
\]

Substituting \(p_q=N_q(e)\) proves (4.2).
An equivalent recursive form is

\[
\boxed{
\Phi_{m,0}=1,\qquad
i\Phi_{m,i}
=
\sum_{r=1}^i
(-1)^{r-1}N_{mr}\Phi_{m,i-r}.
}
\tag{4.5}
\]

For example,

\[
\begin{aligned}
\Phi_{m,1}&=N_m,\\[1mm]
\Phi_{m,2}&=\frac{N_m^2-N_{2m}}2,\\[1mm]
\Phi_{m,3}
&=
\frac{N_m^3-3N_mN_{2m}+2N_{3m}}6.
\end{aligned}
\tag{4.6}
\]


\begin{enumerate}
\item Particularly simple formulas for \(m=2\) and \(m=3\)
\end{enumerate}
5.1 Squares
Taking \(\zeta=-1\) in (2.2) gives

\[
e_i^{[2]}
=
(-1)^i[z^{2i}]E(z)E(-z).
\]

Thus

\[
e_i^{[2]}
=
(-1)^i
\sum_{\substack{a+b=2i\\0\le a,b\le n}}
(-1)^be_ae_b.
\]

Pairing the terms \(a=i-j,\ b=i+j\) yields the especially simple identity

\[
\boxed{
\Sigma_i(X_1^2,\ldots,X_n^2)
=
\Sigma_i^2+
2\sum_{j=1}^{\min(i,n-i)}
(-1)^j\Sigma_{i-j}\Sigma_{i+j},
}
\tag{5.1}
\]

where \(\Sigma_0=1\).
For example,

\[
\begin{aligned}
\Sigma_1(X_1^2,\ldots,X_n^2)
&=\Sigma_1^2-2\Sigma_2,\\
\Sigma_2(X_1^2,\ldots,X_n^2)
&=\Sigma_2^2-2\Sigma_1\Sigma_3+2\Sigma_4,\\
\Sigma_3(X_1^2,\ldots,X_n^2)
&=\Sigma_3^2-2\Sigma_2\Sigma_4
+2\Sigma_1\Sigma_5-2\Sigma_6.
\end{aligned}
\]

Terms whose indices exceed \(n\) are understood to be zero.
5.2 Cubes
Let \(\omega^3=1\), \(\omega\ne1\). Formula (2.1) becomes

\[
\boxed{
\Sigma_i(X_1^3,\ldots,X_n^3)
=
\sum_{\substack{0\le a,b,c\le n\\a+b+c=3i}}
\omega^{b+2c}\Sigma_a\Sigma_b\Sigma_c.
}
\tag{5.2}
\]

Grouping equal monomials gives a formula with visibly integral coefficients:

\[
\boxed{
\begin{aligned}
\Sigma_i(X_1^3,\ldots,X_n^3)
={}&\Sigma_i^3\\
&+3
\sum_{\substack{0\le a,c\le n\\a\ne c,\ 2a+c=3i}}
\Sigma_a^2\Sigma_c\\
&+6
\sum_{\substack{0\le a<b<c\le n\\
a+b+c=3i\\
a\equiv b\equiv c\pmod3}}
\Sigma_a\Sigma_b\Sigma_c\\
&-3
\sum_{\substack{0\le a<b<c\le n\\
a+b+c=3i\\
\{a,b,c\}\bmod3=\{0,1,2\}}}
\Sigma_a\Sigma_b\Sigma_c.
\end{aligned}
}
\tag{5.3}
\]

Indeed, for a multiset \(\{a,a,c\}\) the three root-of-unity weights sum to \(3\). For three distinct indices whose sum is divisible by \(3\), their residues are either all equal, giving weight sum \(6\), or are \(0,1,2\), giving

\[
3\omega+3\omega^2=-3.
\]

For instance,

\[
\Sigma_1(X_1^3,\ldots,X_n^3)
=
\Sigma_1^3-3\Sigma_1\Sigma_2+3\Sigma_3.
\]

For \(n=4\),

\[
\begin{aligned}
\Sigma_2(X_1^3,\ldots,X_4^3)
={}&
\Sigma_2^3
-3\Sigma_1\Sigma_2\Sigma_3
+3\Sigma_1^2\Sigma_4\\
&-3\Sigma_2\Sigma_4+3\Sigma_3^2.
\end{aligned}
\]


\begin{enumerate}
\item Boundary cases and characteristic
\end{enumerate}
The formulas are universal polynomial identities over \(\mathbb Z\). Consequently, they remain valid after specialization in every commutative ring, with no restriction on characteristic.
Several checks follow immediately.
Top elementary symmetric polynomial

\[
\Sigma_n(X_1^m,\ldots,X_n^m)
=
(X_1\cdots X_n)^m
=
\Sigma_n^m.
\]

Thus

\[
\Phi_{m,n}^{(n)}=Y_n^m.
\]

First elementary symmetric polynomial

\[
\Sigma_1(X_1^m,\ldots,X_n^m)
=
X_1^m+\cdots+X_n^m
=
N_m(\Sigma_1,\ldots,\Sigma_n).
\]

One variable
For \(n=1\),

\[
\Phi_{m,1}^{(1)}(Y_1)=Y_1^m.
\]

The omitted case \(m=1\)
The same formulas give

\[
\Phi_{1,i}^{(n)}=Y_i.
\]

Positive characteristic
Let the coefficient ring have characteristic \(p>0\). For \(a\ge0\) and \(q\ge1\),

\[
\Sigma_i(X_1^{p^aq},\ldots,X_n^{p^aq})
=
\left(
\Sigma_i(X_1^q,\ldots,X_n^q)
\right)^{p^a}.
\]

Therefore, as polynomials over \(\mathbb F_p\),

\[
\boxed{
\Phi_{p^aq,i}^{(n)}
\equiv
\left(\Phi_{q,i}^{(n)}\right)^{p^a}
\pmod p.
}
\]

In particular,

\[
\Phi_{p^a,i}^{(n)}
\equiv Y_i^{p^a}\pmod p.
\]

For example, in characteristic \(2\), formula (5.1) reduces to

\[
\Sigma_i(X_1^2,\ldots,X_n^2)=\Sigma_i^2,
\]

as required by Frobenius.
The recurrence (4.5) involves division by \(i\), so it must not be used literally in characteristic dividing \(i\). The companion determinant (2.4), resultant (2.5), and root-of-unity polynomial after integral simplification have no such restriction.

\begin{enumerate}
\item Exact reproducible computation
\end{enumerate}
The following SymPy code uses the division-free companion-matrix formula.
PythonExécuterfrom \_\_future\_\_ import annotations

from itertools import combinations
import sympy as sp


def q831\_polynomials(n: int, m: int):
    """
    Return (e, Phi), where

        Phi[i] = Sigma\_i\(X_1^m, ..., X_n^m\)

    expressed as a polynomial in e = (e1, ..., en).
    """
    if not isinstance(n, int) or not isinstance(m, int) or n < 1 or m < 1:
        raise ValueError("n and m must be positive integers")

    e = tuple(sp.Symbol(f"e\{j\}") for j in range(1, n + 1))
    T = sp.Symbol("T")

\subsection{Companion matrix of}
\subsection{Z\textasciicircum{}n - e1 Z\textasciicircum{}\(n-1\) + e2 Z\textasciicircum{}\(n-2\) - ... + (-1)\textasciicircum{}n en.}
    C = sp.zeros\(n\)

    for r in range(n - 1):
        C[r + 1, r] = 1

    for r in range\(n\):
        j = n - r
        C[r, n - 1] = (-1) ** (j - 1) * e[j - 1]

    characteristic = sp.Poly(
        sp.expand((T * sp.eye\(n\) - C**m).det()),
        T,
    )

    Phi = tuple(
        sp.expand(
            (-1) ** i
\begin{itemize}
\item characteristic.coeff\_monomial(T ** (n - i))
\end{itemize}
        )
        for i in range(n + 1)
    )

    return e, Phi


def elementary(values, k: int):
    """Elementary symmetric polynomial of degree k in values."""
    if k == 0:
        return sp.Integer\(1\)

    return sum(
        (
            sp.prod(values[j] for j in subset)
            for subset in combinations(range\(len(values)\), k)
        ),
        sp.Integer\(0\),
    )


def verify\_q831(n: int, m: int) -> bool:
    """
    Independently substitute the elementary symmetric polynomials
    in x\_1, ..., x\_n and verify every identity symbolically.
    """
    e, Phi = q831\_polynomials(n, m)
    x = tuple(sp.Symbol(f"x\{j\}") for j in range(1, n + 1))

    substitution = \{
        e[j - 1]: elementary(x, j)
        for j in range(1, n + 1)
    \}

    xm = tuple(value**m for value in x)

    return all(
        sp.expand(
            Phi[i].subs\(substitution\) - elementary(xm, i)
        ) == 0
        for i in range(n + 1)
    )


\subsection{Example:}
assert verify\_q831(4, 3)

e, Phi = q831\_polynomials(4, 3)
for i, expression in enumerate\(Phi\):
    print(i, expression)
For \(n=4,m=3\), this returns

\[
\begin{aligned}
\Phi_{3,0}&=1,\\
\Phi_{3,1}&=e_1^3-3e_1e_2+3e_3,\\
\Phi_{3,2}&=
e_2^3-3e_1e_2e_3+3e_1^2e_4
-3e_2e_4+3e_3^2,\\
\Phi_{3,3}&=
e_3^3-3e_2e_3e_4+3e_1e_4^2,\\
\Phi_{3,4}&=e_4^3.
\end{aligned}
\]

Final conclusion
Q831 has a complete affirmative answer. For every \(m\ge2\) and \(1\le i\le n\), the polynomial

\[
\Sigma_i(X_1^m,\ldots,X_n^m)
\]

is given explicitly by the root-of-unity sum (2.1), equivalently by the integral companion determinant/resultant (2.3)--(2.5), or by the finite partition expansion (4.2). All formulas are universal over \(\mathbb Z\), hence valid without exception over every commutative ring and in every characteristic.
